% \iffalse meta-comment
%
%% regulatory-sign.dtx
%% Copyright 2024-2026 E. Nijenhuis
%
% This work may be distributed and/or modified under the
% conditions of the LaTeX Project Public License, either version 1.3c
% of this license or (at your option) any later version.
% The latest version of this license is in
% http://www.latex-project.org/lppl.txt
% and version 1.3c or later is part of all distributions of LaTeX
% version 2005/12/01 or later.
%
% This work has the LPPL maintenance status ‘maintained’.
%
% The Current Maintainer of this work is E. Nijenhuis.
%
% This work consists of the files listed in the meta-comment of
% regulatory-struct.dtx.
%
% \fi
%
% \iffalse
%<*driver>
\ProvidesFile{regulatory-sign.dtx}
%</driver>
%<package>\NeedsTeXFormat{LaTeX2e}
%<package>\ProvidesPackage{regulatory-sign}
%<package>    [2026/09/10 1.0.0 Xerdi's Regulatory Package (Signatures)]
%
%<*driver>
\documentclass[10pt,english]{ltxdoc}
%! suppress = InclusionLoop
\usepackage{regulatory}
\usepackage{tabularx}
\usepackage[english,dutch]{babel}
%% regulatory-preamble.tex
%% Copyright 2024 E. Nijenhuis
%
% This work may be distributed and/or modified under the
% conditions of the LaTeX Project Public License, either version 1.3c
% of this license or (at your option) any later version.
% The latest version of this license is in
% http://www.latex-project.org/lppl.txt
% and version 1.3c or later is part of all distributions of LaTeX
% version 2005/12/01 or later.
%
% This work has the LPPL maintenance status ‘maintained’.
%
% The Current Maintainer of this work is E. Nijenhuis.
%
% This work consists of the files regulatory.ins,
% regulatory-struct.dtx, regulatory-ref.dtx,
% regulatory-defs.dtx, regulatory-attachments.dtx,
% regulatory-md.dtx,
% regulatory-html.dtx,
% regulatory-english.def, regulatory-dutch.def,
% regulatory.tex, regulatory-preamble.tex,
% regulatory-nl.tex, regulatory-en.tex,
% example1.bib, example2.bib,
% example1-nl.tex, example2-nl.tex,
% example1-en.tex, example2-en.tex,
% md-example.tex, example.md,
% and the derived files regulatory.sty, regulatory-struct.sty,
% regulatory-defs.sty, regulatory-ref.sty,
% regulatory-attachments.sty, regulatory-md.sty
% and regulatory.4ht.
\usepackage{gitinfo-lua}
\usepackage{listings}
\usepackage{adjustbox}
\usepackage{tikz}
\usetikzlibrary{arrows.meta}

\usepackage{fontspec}
\usepackage{textcomp}

\def\projecturl{https://github.com/Xerdi/regulatory}

\definecolor{greycolor}{HTML}{BFBFBF}
\definecolor{primary}{HTML}{174A67}
\definecolor{darkcolor}{HTML}{091D28}
\definecolor{greencolor}{HTML}{176734}
\colorlet{artlabel}{primary}
\colorlet{green}{greencolor}

\lstdefinestyle{tex}{%
    language=[LaTeX]TeX,
    keywordsprefix={\\},
    alsoletter={\\},
    morekeywords={article,para,textfill,masterdocument,refdocument,loadglsdefs,printdefs,setmiddleconjunction,setlastconjunction,setrangeconjunction,setconjunction,markdownInput,glssetwidest,describe}
}

\lstdefinestyle{bash}{%
    language=bash,
    morekeywords={pdflatex,lualatex,latexmk,makeglossaries,bib2gls}
}

\lstdefinelanguage{Markdown}[LaTeX]{TeX}{%
    sensitive=false,
    morecomment=[l][\bfseries\ttfamily\color{green}]{\#},
    morecomment=[s][\bfseries\ttfamily\color{purple}]{:\ \ }{\ },
    morestring=[s][\color{black}]{\{}{\}},
    morekeywords={article,para,textfill,masterdocument,refdocument,loadglsdefs,printdefs,setmiddleconjunction,setlastconjunction,setrangeconjunction,setconjunction,markdownInput,glssetwidest,describe},
    keywordsprefix={\\},
    alsoletter={\\},
}

\lstdefinestyle{md}{%
    language=Markdown
}

\lstdefinelanguage{BibTeX}{%
    keywords={%
        article,book,collectedbook,conference,electronic,entry,ieeetranbstctl,%
        inbook,incollectedbook,incollection,injournal,inproceedings,%
        manual,mastersthesis,misc,patent,periodical,phdthesis,preamble,%
        proceedings,standard,string,techreport,unpublished%
    },
    keywordsprefix={@},
    comment=[l][\itshape]{@comment},
    sensitive=false,
}

\lstdefinestyle{bib}{%
    language=BibTeX,
    morestring=[s][\bfseries\ttfamily\color{purple}]{\ \ \ \ }{\ },
    morecomment=[n][\bfseries\ttfamily\color{green}]{\ \{}{\}},
}

\lstset{
    title=\lstname,
    basicstyle=\ttfamily,
    numberstyle=\ttfamily,
    numbersep=5pt,
    rulecolor=\color{greycolor},
    stringstyle=\color{pink},
    keywordstyle=\color{primary},
    commentstyle=\itshape\color{darkcolor},
    breaklines=true,
    captionpos=b,
    tabsize=4,
    showtabs=true
}

\newcommand\package{\texttt}
\newcommand\file{\texttt}
\newcommand\option{\texttt}

% Reads test/conformance.tex, which test/conformance.sh generates and which is
% kept in the repository: veraPDF is not available everywhere this manual is
% built, so the verdicts travel with the sources.
\newcommand\conformanceversion{?}
\newcommand\conformanceimage{?}
% A digest carries no break points of its own and does not fit a line, so this
% offers one after every character. The argument is expanded first, otherwise the
% loop sees one macro instead of the characters it stands for.
\newcommand\anywherebreakend{}
\def\anywherebreakloop#1{%
    \ifx#1\anywherebreakend\else#1\allowbreak\expandafter\anywherebreakloop\fi}
\newcommand\anywherebreak[1]{\anywherebreakloop#1\anywherebreakend}
\newcommand\conformanceimagetext{%
    \expandafter\anywherebreak\expandafter{\conformanceimage}}

% \ignorespaces, because this macro typesets nothing while its line in the
% generated file still ends in a newline, and that space would otherwise open
% the first cell of the table.
\newcommand\conformancevalidator[2]{%
    \gdef\conformanceversion{#1}\gdef\conformanceimage{#2}\ignorespaces}
\newcommand\conformancerow[6]{%
    \texttt{#1} & \texttt{#2} & \texttt{#3} & \texttt{#4} & \texttt{#5}%
    \ifthenelse{\equal{#6}{}}{}{ \footnotesize(#6)} \\}

% Points at the resulting PDF of an example, which is attached to this manual.
% Its argument is the label of the artifact the example was declared with, and it
% falls back to plain text whenever the PDF isn't there.
\newcommand\exampleresult[1]{%
    \par\noindent
    \translation{The result of this example is}{Het resultaat van dit voorbeeld is}
    \documentattachment{#1}{\file{\artifactfilename{#1}}}.\par}

\def\packagedate{\gitdate}
\def\packageversion{\gitversion}

%\def\cmd{\lstinline[style=tex]}
\def\cmditem#1{\item[\ttfamily\textbackslash{}#1]}

\newcommand\Vtextvisiblespace[1][.3em]{%
  \mbox{\kern.06em\vrule height.3ex}%
  \vbox{\hrule width#1}%
  \hbox{\vrule height.3ex}}

% Package zref-clever loads its language files lazily, which would otherwise
% happen inside a \DocInput region, where the percent sign is ignored and the
% comments of those files end up being parsed as keys. Touching every language
% of this manual schedules them at the start of the document instead, while the
% percent sign still is a comment character.
\zcsetup{lang=english}
\zcsetup{lang=dutch}
\zcsetup{lang=current}

\newcommand\translation[2]{#1}
\begin{document}
    \selectlanguage{english}
    \DocInput{regulatory-sign.dtx}
\end{document}
%</driver>
% \fi
%
% \subsection{\texorpdfstring{\package{regulatory-sign}}{regulatory-sign}}
% \setcounter{CodelineNo}{0}
%
% This module comes from \package{xdp-sign} of Xerdi's Documentation Project, where it was written for the
% agreement class of that bundle. A signature belongs to the same kind of document as the rest of this
% package, and a document that needs one should not have to reach for a package that is in no
% distribution, so it is maintained here.
%
% \subsubsection{\translation{Prerequisites}{Vereisten}}
%
% A signature field is a widget annotation, which needs \package{hyperref} for the form machinery it
% belongs to. It is loaded by \package{regulatory-struct}, which is stated here for the document that
% loads this module on its own. Options for \package{hyperref} have to be passed before either is loaded,
% for instance with \cmd{\PassOptionsToPackage}\marg{hidelinks}\marg{hyperref}.
% \iffalse
%<*package>
% \fi
%    \begin{macrocode}
\RequirePackage{regulatory-struct}
\RequirePackage{hyperref}
%    \end{macrocode}
%
% \subsubsection{\translation{The annotation}{De annotatie}}
%
% \begin{macro}{\regulatory@sign@annot}
% Draws a box of \meta{width} by \meta{height} with a rule at the top and at the bottom, and puts a PDF
% annotation of exactly that rectangle over it, so that a viewer offers the field where the box is.
%
% There are two ways to place one, and which is available depends on the document rather than on the
% engine. With PDF management active\,---\,which is what \cmd{\DocumentMetadata} turns on, and what
% \package{regulatory-attachments} needs anyway\,---\,the annotation goes through
% \verb|\pdfannot_box:nnnn| of the kernel, which knows about the page it lands on and about the standards
% the document declares. Without it, the engine primitive is the only route: \verb|\pdfextension annot|
% under Lua\TeX{} and current pdf\TeX{}, \cmd{\pdfannot} under an older one. Measured: the kernel command
% exists only under \cmd{\DocumentMetadata}, so the test is for the command and not for the engine.
%
% A document that has neither gets a warning and a drawn box without a field, which is a signature line to
% write on but not one to sign in a viewer. That is worth saying out loud: a missing annotation is
% invisible on paper and the difference only shows when someone tries to sign.
%    \begin{macrocode}
\newcommand*\regulatory@sign@annot[3]{%
    \begingroup
    \setlength{\dimen0}{#1}%
    \setlength{\dimen1}{#2}%
    \leavevmode
    \hbox to \dimen0{%
        \vbox to \dimen1{%
            \hrule height0.4pt\relax
            \vfill
            \hrule height0.4pt\relax
        }\hss
    }%
    \kern-\dimen0
    \raisebox{0pt}[0pt][0pt]{%
        \hbox to 0pt{%
            \hskip0pt
            \regulatory@sign@place{#3}%
            \hss
        }%
    }%
    \kern\dimen0
    \endgroup
}
%    \end{macrocode}
% \end{macro}
%
% \begin{macro}{\regulatory@sign@place}
% Emits the annotation itself, by whichever of the three routes this document has.
%    \begin{macrocode}
\ExplSyntaxOn
\cs_new_protected:Npn \regulatory@sign@place #1
{
    \cs_if_exist:NTF \pdfannot_box:nnnn
        { \pdfannot_box:nnnn { \dimen0 } { \dimen1 } { 0pt } { #1 } }
        {
            \cs_if_exist:NTF \pdfextension
                { \tex_pdfextension:D annot ~ width ~ \dimen0 ~
                  height ~ \dimen1 ~ depth ~ 0pt ~ { #1 } }
                {
                    \cs_if_exist:NTF \pdfannot
                        { \pdfannot ~ width ~ \dimen0 ~
                          height ~ \dimen1 ~ depth ~ 0pt ~ { #1 } }
                        {
                            \msg_warning:nn { regulatory-sign } { no-annotation }
                        }
                }
        }
}

\msg_new:nnn { regulatory-sign } { no-annotation }
    {
        This~engine~offers~no~way~to~place~an~annotation,~so~the~signature~\\
        field~is~a~drawn~box~and~nothing~more.~Declare~\iow_char:N\\DocumentMetadata~
        to~let~the~kernel~place~it.
    }
\ExplSyntaxOff
%    \end{macrocode}
% \end{macro}
%
% \subsubsection{\translation{Signature fields}{Handtekeningvelden}}
%
% \begin{macro}{\regulatory@sign@appearance}
% An annotation has to say what it looks like. ISO 19005 puts it as a requirement rather than a nicety:
% rule 6.3.3-1 of both PDF/A-2 and PDF/A-3 asks every annotation that is not a \texttt{Popup}, a
% \texttt{Link} or of zero size to carry an appearance dictionary, and veraPDF says of a field without one
% \enquote{An annotation does not contain an appearance dictionary}. A signature field is a
% \texttt{Widget}, so it is asked as well, and a document with one used to lose its conformance over it.
%
% What it points at is deliberately empty. The line to sign on is drawn on the page by
% \cmd{\regulatory@sign@annot}, not by the annotation, so an appearance that drew anything would draw it
% twice; and a signing tool replaces the normal appearance of the field with its own the moment the field
% is signed. One object serves every field: the bounding box of an appearance is mapped onto the rectangle
% of the annotation, so a stream that draws nothing needs no size of its own.
%
% Only where the PDF management of \LaTeX{} is active, which a document that cares about PDF/A has
% declared anyway, since the standards need it for much more than this. The test is for the management
% and not for the command: under \package{tex4ht} the object commands exist and their backend does not,
% so asking whether they are defined answers yes and then fails on the first call.
%    \begin{macrocode}
\ExplSyntaxOn
\cs_new:Npn \regulatory@sign@appearance { }
\bool_lazy_and:nnT
  { \cs_if_exist_p:N \pdfmanagement_if_active_p: }
  { \cs_if_exist_p:N \pdf_object_new:n }
  {
    \pdfmanagement_if_active:T
      {
    \pdf_object_new:n { regulatory/sign/blank }
    \AddToHook { shipout/firstpage }
      {
        \pdf_object_write:nnn { regulatory/sign/blank } { stream }
          { { /Type /XObject /Subtype /Form /BBox [0~0~1~1] } { } }
      }
    \cs_gset:Npn \regulatory@sign@appearance
      { /AP~<<~/N~\pdf_object_ref:n { regulatory/sign/blank }~>> }
      }
  }
\ExplSyntaxOff
%    \end{macrocode}
% \end{macro}
%
% \begin{macro}{\signaturefield}
% Draws a signature field of \meta{width} by \meta{height} under the name \meta{name}, which is the name a
% viewer stores the signature under and the one a signing tool asks for. The optional \meta{tooltip} is
% what a reader is told the field is for, which is also what a screen reader announces, so it is worth
% giving.
%    \begin{macrocode}
\newcommand\signaturefield[4][]{%
    \def\regulatory@sign@tooltip{#1}%
    \regulatory@sign@annot{#3}{#4}{%
        /Subtype /Widget
        /FT /Sig
        /T (#2)
        \ifx\regulatory@sign@tooltip\@empty\else /TU (#1)\fi
        /BS << /W 0 >>
        /F 4
        /DA (/Helv 0 Tf 0 g)
        \regulatory@sign@appearance
    }%
}
%    \end{macrocode}
% \end{macro}
%
% \begin{macro}{\SignatureField}
% The name this command had in \package{xdp-sign}, kept so that a document written for that package keeps
% working when it is moved over. New documents use the lowercase name, which is the one the rest of this
% bundle is written in.
%    \begin{macrocode}
\let\SignatureField\signaturefield
\let\SignatureAnnotation\regulatory@sign@annot
%    \end{macrocode}
% \end{macro}
% \iffalse
%</package>
% \fi
%
% \Finale
%
